\section{Sufficient conditions weaker than linear independence}\label{sec:sufficient}

In this section GRH$(\chi)$ and NCZ$(\chi)$ are assumed throughout. Each hypothesis below is a condition on the set $\Gamma$ (with multiplicities); none is known to hold for any $L$-function, and each is implied by LI (with simple zeros where stated). For $X\subset\RR$, $\spanQ X$ denotes the $\QQ$-span. In every case the conclusion ``no atoms'' gives, by \eqref{eq:portmanteau} and \eqref{eq:sandwich} exactly as in Theorem~\ref{thm:jump}(d): for every $m>-\frac12$ the logarithmic densities of $\{D_m>0\}$, $\{D_m\ge0\}$ and $\{x:D_m(x)>c\,x^M/\log x\}$ (every $c$) exist, and near-ties have logarithmic density $0$. In particular, for $\chi_{-4}$, the logarithmic density of $\{x:\pi(x;4,3)>\pi(x;4,1)\}$ exists.

\subsection{The tail does not span: TS, FD and tame directions}\label{sec:TS}

\begin{hypothesis}[TS]
There is $T>0$ with $\spanQ{\Gamma_{>T}}\ne\spanQ\Gamma$, i.e.\ some ordinate is not a rational linear combination of ordinates exceeding $T$. Equivalently, there is a $\QQ$-linear form $\varphi\ne0$ on $\spanQ\Gamma$ with $\varphi(\gamma)=0$ for all but finitely many $\gamma\in\Gamma$.
\end{hypothesis}

\begin{theorem}[GRH, NCZ, TS; \rd{1 read}]\label{thm:TS}
For every $M>0$ the measure $\nu_M$ is absolutely continuous.
\end{theorem}

\begin{proof}
Let $\varphi$ be as in TS and $\Lambda=\{\gamma\in\Gamma:\varphi(\gamma)\ne0\}$ (finite, nonempty); after scaling, $\varphi(G)\subset\ZZ$. Put $j(s)=(e^{i\varphi(\gamma)s})_{\gamma\in\Gamma}$. If $\sum n_\gamma\gamma=0$ then $\prod e^{i\varphi(\gamma)sn_\gamma}=1$, so $j(s)\in\cA$ by (I2), and $j:\RR\to\cA$ is a $2\pi$-periodic one-parameter subgroup with $e_g(j(s))=e^{i\varphi(g)s}$. Write $F_M=F'+P$ with $P=\re\sum_{\gamma\in\Lambda}m_\gamma r_\gamma e_\gamma$. The Fourier coefficients of $F'$ lie in $\ker\varphi$, so for every $s$, $F'(\theta j(s))=F'(\theta)$ for a.e.\ $\theta$; by Fubini there is a null set $N_1$ such that for $\theta\notin N_1$ this holds for a.e.\ $s\in[0,2\pi]$. For every $\theta$, $p_\theta(s):=P(\theta j(s))=\sum_{k\ne0}C_k(\theta)e^{iks}$ with
\[
C_k(\theta)=\tfrac12\sum_{\gamma\in\Lambda,\ \varphi(\gamma)=k}m_\gamma r_\gamma\theta_\gamma+\tfrac12\sum_{\gamma\in\Lambda,\ \varphi(\gamma)=-k}m_\gamma\bar r_\gamma\bar\theta_\gamma .
\]
Choose $k_0\ne0$ with $\{\gamma\in\Lambda:\varphi(\gamma)=\pm k_0\}\ne\emptyset$. Then $C_{k_0}$ is a trigonometric polynomial in pairwise distinct characters with nonzero coefficients, so its zero set $N_2$ is null (Lemma~\ref{lem:trigpoly}); for $\theta\notin N_2$, $p_\theta$ is non-constant and real-analytic. For a Lebesgue-null Borel $B$, by invariance and Fubini,
\[
\nu_M(B)=\int_\cA\frac1{2\pi}\mathrm{Leb}\{s\in[0,2\pi]:F_M(\theta j(s))\in B\}\,d\hA(\theta),
\]
and for $\theta\notin N_1\cup N_2$ the inner set equals $\{s:F'(\theta)+p_\theta(s)\in B\}$ up to a null set, which is null by Lemma~\ref{lem:analytic}.
\end{proof}

\begin{corollary}[what a non-absolutely-continuous law forces; GRH, NCZ; \rd{1 read}]\label{cor:TSatom}
If for some $M>0$ the measure $\nu_M$ is not absolutely continuous (in particular if it has an atom), then for every $T$ every ordinate of $L(s,\chi)$ is a rational linear combination of ordinates exceeding $T$.
\end{corollary}

\begin{example}[\rd{1 read}]\label{ex:TSweaker}
TS is strictly weaker than the hypothesis of Devin \cite[Theorem~2.2]{Devin} (a finite nonempty $\Lambda'$ with $\spanQ{\Lambda'}\cap\spanQ{\Gamma\setminus\Lambda'}=\{0\}$), as a condition on a set of frequencies. Let $b_0,b_1,\ldots>0$ be $\QQ$-independent with $b_i\to\infty$, and $X=\{b_0,b_0+b_1\}\cup\{b_i+b_j:1\le i<j\}$. The dual form $b_0^*$ is nonzero exactly on $b_0,b_0+b_1$, so TS holds. But for a finite $\Lambda'$ every $b_i$ ($i\ge1$) lies in $\spanQ{X\setminus\Lambda'}$ ($2b_i=(b_i+b_k)+(b_i+b_l)-(b_k+b_l)$ for large $k<l$), so the intersection is never $\{0\}$.
\end{example}

\begin{hypothesis}[FD, finite-level directions]
There are $N_0\ge1$ and $\eps\in(0,\frac12)$ such that for arbitrarily large $T$ there exist $\Lambda_T\subset\Gamma_T$ with $\#\Lambda_T\le N_0$ and $\gamma_T\in\Lambda_T$ with $\gamma_T\le T^{1/2-\eps}$ and $\gamma_T\notin\spanQ{\Gamma_T\setminus\Lambda_T}$.
\end{hypothesis}

TS implies FD (take the fixed set $\Lambda=\{\gamma:\varphi(\gamma)\ne0\}$).

\begin{theorem}[GRH, NCZ, FD; \rd{1 read}]\label{thm:FD}
For every $M>0$ the measure $\nu_M$ has no atoms. If FD holds for all $T\ge T_0$, then $\sup_c\nu_M([c-\ell,c+\ell])\ll\ell^\beta$ as $\ell\to0$, with $\beta=\eps/(4N_0(1-\eps))$.
\end{theorem}

\begin{proof}
Fix $M$ and $T$ as in FD; there is a homomorphism $\varphi:G_T\to\ZZ$ ($G_T$ generated by $\Gamma_T$) vanishing on $\Gamma_T\setminus\Lambda$ and nonzero at $\gamma_T$, where $\Lambda=\Lambda_T$, $n=\#\Lambda\le N_0$. Let $\cA_T$ be the closure of $\{(e^{i\gamma y})_{\gamma\in\Gamma_T}\}$ and $\pi_T:\cA\to\cA_T$ the restriction (onto, measure preserving). By (I7),
\[
\|F_M-F_{M,T}\circ\pi_T\|_2^2=V_M^{>T}\le C_1M^2\log^2(T+3)/T.
\]
As in Theorem~\ref{thm:TS}, $j(s)=(e^{i\varphi(\gamma)s})_{\gamma\in\Gamma_T}$ is a $2\pi$-periodic subgroup of $\cA_T$ and $F_{M,T}(\theta j(s))=F'(\theta)+\sum_{k\in K}C_k(\theta)e^{iks}$, $\#K\le2n$, $0\notin K$. Put $k_*=\varphi(\gamma_T)$ and $a=\frac12m_{\gamma_T}|r_{\gamma_T}|\ge2M/(M+T^{1/2-\eps})$. Fix $c$, $\ell>0$, $\theta$, let $q(s)=F'(\theta)-c+\sum_KC_k(\theta)e^{iks}$ and $E_\theta=\{s\in[0,2\pi]:|q(s)|\le\ell\}$. Since $\sup|q|\ge|C_{k_*}(\theta)|$, (I6) gives $|E_\theta|/2\pi\le A_N(\ell/|C_{k_*}(\theta)|)^{1/(2N_0)}$ when $\ell\le|C_{k_*}(\theta)|$. Next, $\hA\{|C_{k_*}|\le\eta a\}\le A_N\eta^{1/N_0}$ for $0<\eta<1$: if $C_{k_*}$ has one term the set is empty; otherwise choose $\psi:G_T\to\ZZ$ injective on the finitely many frequencies of $C_{k_*}$, move along $j_\psi$, and apply (I6) to a trigonometric polynomial with at most $N_0$ distinct integer frequencies, one coefficient of modulus $a$. By invariance and Fubini,
\[
\nu_M([c-\tfrac\ell2,c+\tfrac\ell2])\le A_N\eta^{1/N_0}+A_N\Bigl(\frac\ell{\eta a}\Bigr)^{1/(2N_0)}+\frac{4C_1M^2\log^2(T+3)}{T\ell^2}\qquad(\ell\le\eta a).
\]
Take $\ell=T^{-1/2+\eps/2}$ and $\eta=T^{-\eps/4}$; then $\ell/(\eta a)\le T^{-\eps/4}(M+1)/(2M)$, and the right side is $\ll T^{-\eps/(8N_0)}+T^{-\eps}\log^2(T+3)$, uniformly in $c$. Let $T\to\infty$ along FD. If every $T\ge T_0$ is allowed, take $T$ with $T^{-1/2+\eps/2}=2\ell_1$.
\end{proof}

\emph{Comparison with the literature} (statements read in the arXiv versions; \file{density/LITERATURE.md}). One self-sufficient ordinate ($\gamma\notin\spanQ{\Gamma\setminus\{\gamma\}}$; Devin \cite[Theorem~5.4(2)]{Devin}, Martin--Ng \cite[Definition~1.3]{MN}) implies Devin's direct-sum hypothesis \cite[Theorem~2.2]{Devin} (also Bailleul \cite[Theorem~1.7]{Bailleul}, Devin \cite[Corollary~2]{Devin2}), which implies TS, which implies FD. Devin's Theorem~5.4(1) (for every $T\ge T_0$ some $\gamma_T\le T^{1/2-\eps}$ not in $\spanQ{\Gamma_T\setminus\{\gamma_T\}}$) is FD with $N_0=1$ and ``every $T\ge T_0$''. TS does not imply it (Example~\ref{ex:TSweaker}) and it does not imply TS (an example supplied by the first referee: \file{density/DENSITY.md}, Example~3.3a; it was rewritten after the second read, following that read's suggestion, and not re-read). So FD is implied by each of these hypotheses and strictly weaker than each. Theorem~\ref{thm:TS} gives absolute continuity under TS, whereas \cite[Theorem~2.2]{Devin} gives continuity under the direct-sum hypothesis. The gain is modest.

\begin{definition}
A homomorphism $\varphi:G\to\RR$ is \emph{tame} if $|\varphi(\gamma)|\le C_\varphi\log(\gamma+3)$ for all $\gamma\in\Gamma$. For $v\in\RR$ let $C^\varphi_v=\sum_{g\in\pm\Gamma,\ \varphi(g)=v}m_g\frac{2M}{M+ig}e_g\in L^2(\cA)$ (with $m_{-\gamma}=m_\gamma$), and $j_\varphi(s)=(e^{i\varphi(\gamma)s})_\gamma$, a one-parameter subgroup of $\cA$ by (I2).
\end{definition}

\begin{theorem}[tame directions; GRH, NCZ; \rd{1 read}]\label{thm:tame}
Let $\varphi$ be tame. \textup{(a)} For a.e.\ $\theta$ there is a real-analytic $h_\theta$ with $F_M(\theta j_\varphi(s))=h_\theta(s)$ for a.e.\ $s$. \textup{(b)} If $S=\{F_M=c\}$ has positive measure, then $S$ is invariant under $j_\varphi(\RR)$ up to null sets, $C^\varphi_v=0$ a.e.\ on $S$ for every $v\ne0$, and $1+C^\varphi_0=c$ a.e.\ on $S$. \textup{(c)} If for some $v\ne0$ the set $\{g\in\pm\Gamma:\varphi(g)=v\}$ is finite and nonempty, then $\nu_M$ is absolutely continuous.
\end{theorem}

\begin{proof}[Proof (outline; complete proof in \file{density/DENSITY.md}, Theorem~3.5)]
(a) With $c_0=1/(4C_\varphi)$, the partial sums $F^z_N=1+\sum_{|g|\le N}m_g\frac{2M}{M+ig}e^{i\varphi(g)z}\theta_g$ are entire in $z$, and $|e^{i\varphi(g)z}|\le(|g|+3)^{1/4}$ on $|\im z|\le c_0$; by orthonormality and (I7) they are $L^2$-Cauchy uniformly on the strip, and the sub-mean-value property gives a.e.\ locally uniform convergence along a subsequence to a holomorphic $H_\theta$, with $F^s_N(\theta)=F^0_N(\theta j_\varphi(s))$. (b) For a.e.\ $\theta\in S$, $h_\theta=c$ on a set of positive measure, hence everywhere; so $S$ is $j_\varphi$-invariant, and the projections $P_v=\lim\frac1R\int_0^Re^{-ivs}U_s\,ds$ (mean ergodic theorem along $j_\varphi$), which commute with multiplication by $\mathbf 1_S$, applied to $\mathbf 1_SF_M=c\mathbf 1_S$ give the claims. (c) If $\nu_M(B)>0$ for a null Borel $B$, the set where $h_\theta$ is constant has positive measure, and the argument of (b) gives $C^\varphi_v=0$ on it, contradicting Lemma~\ref{lem:trigpoly} for the nonzero trigonometric polynomial $C^\varphi_v$.
\end{proof}

\begin{corollary}[necessary conditions for an atom; GRH, NCZ; \rd{1 read}]\label{cor:tame}
If $\nu_M(\{c\})>0$ and $S=\{F_M=c\}$, then for every tame $\varphi$: $\mathbf 1_S$ is $j_\varphi(\RR)$-invariant; every fibre $\{g\in\pm\Gamma:\varphi(g)=v\}$, $v\ne0$, is empty or infinite; and if $c\ne1$ the fibre over $0$ is infinite.
\end{corollary}

\subsection{The tail does not generate: TS\texorpdfstring{$_\ZZ$}{\_Z}}\label{sec:TSZ}

\begin{proposition}[finite-index criterion; GRH, NCZ; \rd{2 reads}]\label{prop:index}
Let $H\subset G$ be a subgroup of finite index $N$ such that $\Lambda=\Gamma\setminus H$ is finite. Then $\nu_M(\{c\})\le1/N$ for every $M>0$ and $c\in\RR$.
\end{proposition}

\begin{proof}
Let $K=H^\perp=\{\theta\in\cA:e_h(\theta)=1\ \forall h\in H\}$, a subgroup of order $N$ (dual of $G/H$ by (I2)). Since $G=H+\langle\Lambda\rangle$ and the characters separate points, for $k\ne k'$ in $K$ some $\gamma\in\Lambda$ has $e_\gamma(k)\ne e_\gamma(k')$. Write $F_M=F_H+F_\Lambda$ with $F_\Lambda=\re\sum_{\gamma\in\Lambda}m_\gamma r_\gamma e_\gamma$; the coefficients of $F_H$ lie in $H$, so $F_H(\cdot k)=F_H$ a.e., and after replacing $F_H$ by its average over $K$ this holds everywhere. Let $E=\{F_M=c\}$ and $S(\theta)=\{k\in K:\theta k\in E\}$; then $\int|S(\theta)|\,d\hA=N\nu_M(\{c\})$. If $k\ne k'$ both lie in $S(\theta)$, then $F_\Lambda(\theta k)=F_\Lambda(\theta k')$, so $\theta$ is a zero of the trigonometric polynomial $\re\sum_{\gamma\in\Lambda}m_\gamma r_\gamma(e_\gamma(k)-e_\gamma(k'))e_\gamma$, which has a nonzero coefficient; by Lemma~\ref{lem:trigpoly} this happens on a null set. So $|S(\theta)|\le1$ a.e.
\end{proof}

\begin{hypothesis}[TS$_\ZZ$]
$[G:\langle\Gamma_{>T}\rangle]\to\infty$ as $T\to\infty$ (the index is nondecreasing, and counts as $\infty$ when infinite).
\end{hypothesis}

\begin{theorem}[GRH, NCZ, TS$_\ZZ$; \rd{2 reads}]\label{thm:TSZ}
For every $M>0$ the measure $\nu_M$ has no atoms.
\end{theorem}

\begin{proof}
If some index $[G:\langle\Gamma_{>T}\rangle]$ is infinite, TS holds: otherwise every $\lambda\in\Gamma_T$ has a positive multiple in $\langle\Gamma_{>T}\rangle$ (as $\spanQ{\Gamma_{>T}}=\spanQ\Gamma$), and $G/\langle\Gamma_{>T}\rangle$ is generated by finitely many elements of finite order. Then Theorem~\ref{thm:TS} applies. Otherwise Proposition~\ref{prop:index} with $H=\langle\Gamma_{>T}\rangle$ gives $\nu_M(\{c\})\le1/[G:H]\to0$.
\end{proof}

\begin{corollary}[what an atom forces; GRH, NCZ; \rd{2 reads}]\label{cor:TSZatom}
If $\nu_M(\{c\})=\alpha>0$, then $\langle\Gamma_{>T}\rangle$ is eventually constant, equal to some $H_\infty$ with $[G:H_\infty]\le1/\alpha$, and for every $T$ each ordinate $\gamma$ has a multiple $n\gamma$, $1\le n\le1/\alpha$, that is an integer linear combination of ordinates exceeding $T$.
\end{corollary}

TS implies TS$_\ZZ$, and TS$_\ZZ$ is strictly weaker: $X=\{2^k\alpha:k\ge0\}$ has $\langle X_{>T}\rangle=2^{k(T)}\ZZ\alpha$, of index $\to\infty$, while TS fails. TS$_\ZZ$ and FD are incomparable ($X$ fails FD; the set of \file{density/DENSITY.md}, Example~3.3a, satisfies FD and fails TS$_\ZZ$). Devin's parity hypothesis \cite[Theorem~2.3]{Devin} implies neither ($X=\{(2k+1)\alpha:k\ge0\}$). An atom of mass $\alpha$ makes $G/H_\infty$ a finite group of order $\le1/\alpha$; it does not force it to be trivial, and Corollary~\ref{cor:TSZatom} gives no contradiction with the certified bound of Theorem~\ref{thm:cert}.

\subsection{A low vertex: hypothesis V}\label{sec:V}
For a trigonometric polynomial $P=\sum_{g\in S}c_ge_g$ on the dual of a torsion-free group ($c_g\ne0$), a point $g_0\in S$ is a \emph{vertex} if $g_0\notin\mathrm{conv}(S\setminus\{g_0\})$ in $G\otimes\RR$; equivalently some homomorphism $\varphi:\langle S\rangle\to\ZZ$ has $\varphi(g_0)>\varphi(g)$ for $g\in S\setminus\{g_0\}$. Write $\mathfrak m(f)=\int\log|f|$ (the logarithmic Mahler measure).

\begin{lemma}[Jensen--Newton; classical; \rd{1 read}]\label{lem:JN}
If $g_0$ is a vertex of $\mathrm{conv}(S)$, then $\mathfrak m(P)\ge\log|c_{g_0}|$.
\end{lemma}

\begin{proof}
Reduce to the torus dual to $\langle S\rangle$ and choose $\varphi$ as above. For each $\theta$, $s\mapsto P(\theta j(s))$ equals $e^{ik_-s}p_\theta(e^{is})$ for a polynomial $p_\theta$ of degree $k_+-k_-$ whose leading coefficient is $c_{g_0}e_{g_0}(\theta)$ ($k_\pm$ the extreme values of $\varphi$ on $S$). By Jensen's formula, $\frac1{2\pi}\int_0^{2\pi}\log|p_\theta(e^{is})|ds\ge\log|c_{g_0}|$ for every $\theta$; integrate over $\theta$ (Tonelli, separately for $\log^\pm$). (The lemma and this proof are classical; for $\gamma_0$ maximal in a group order, part (a) of Theorem~\ref{thm:V} below also follows from Jensen's inequality for $H^1$ of a compact group with ordered dual, Helson--Lowdenslager \cite{HL}.)
\end{proof}

\begin{lemma}[small ball from the Mahler measure; \rd{1 read}]\label{lem:smallball}
On a probability space, for $f\in L^2$ with $\mathfrak m(f)>-\infty$ and $0<\delta<1$: $\mu\{|f|<\delta\}\log(1/\delta)\le-\mathfrak m(f)+\frac12\log^+\int|f|^2+\frac1{2e}$.
\end{lemma}

\begin{proof}
With $B=\{|f|<\delta\}$, $\eta=\mu(B)$: $\int_B\log|f|\le\eta\log\delta$, and Jensen's inequality on $B^c$ gives $\int_{B^c}\log|f|\le\frac12\log^+\int|f|^2+\frac{1-\eta}2\log\frac1{1-\eta}$; and $x\log(1/x)\le1/e$.
\end{proof}

\begin{hypothesis}[V$(\gamma_0)$]
For $\gamma_0\in\Gamma$: for every $T\ge\gamma_0$, $\gamma_0$ is a vertex of $\mathrm{conv}(\pm\Gamma_T)$. Equivalently \textup{(\file{density3/MAHLER.md}, Proposition~3.1)}: every nontrivial integer relation $\sum k_\gamma\gamma=0$ among the distinct ordinates satisfies $|k_{\gamma_0}|<\sum_{\gamma\ne\gamma_0}|k_\gamma|$; equivalently, some total group order on $G$ has $\gamma_0=\max(\pm\Gamma)$.
\end{hypothesis}

\begin{theorem}[GRH, NCZ, V$(\gamma_0)$; \rd{1 read}]\label{thm:V}
For every $M>0$: \textup{(a)} $\int\log|x-c|\,d\nu_M(x)\ge\log(2Mm_{\gamma_0}/|M+i\gamma_0|)$ for every $c\in\CC$; \textup{(b)} for $c\in\RR$, $0<\eps<1$,
\[
\nu_M((c-\eps,c+\eps))\le\frac{\log\frac{|M+i\gamma_0|}{2Mm_{\gamma_0}}+\frac12\log^+((1-c)^2+V_M)+\frac1{2e}}{\log(1/\eps)};
\]
\textup{(c)} $\nu_M$ is absolutely continuous.
\end{theorem}

\begin{proof}
(a) For $T\ge\gamma_0$, the support of $F_{M,T}-c$ is $\pm\Gamma_T$ together with $0$, the midpoint of $\pm\gamma_0$, so $\gamma_0$ is a vertex and Lemma~\ref{lem:JN} gives $\mathfrak m(F_{M,T}-c)\ge\log(2Mm_{\gamma_0}/|M+i\gamma_0|)$. As $F_{M,T}\circ\pi_T\to F_M$ in $L^2$, reverse Fatou applied to $|f_T|-\log|f_T|\ge0$ along an a.e.-convergent subsequence gives the bound for $F_M$. (b) Lemma~\ref{lem:smallball}. (c) A compactness argument (inverse limit of the nonempty compact sets $\{\varphi\in\mathrm{Hom}(G_T,\RR):\varphi(\gamma_0)=1,\ |\varphi(\gamma)|\le1\ \forall\gamma\in\Gamma_T\}$) gives a homomorphism $\varphi:G\to\RR$ with $\varphi(\gamma_0)=1$ and $|\varphi|\le1$ on $\Gamma$; it is tame. As in (a), $\mathfrak m(C^\varphi_1)>-\infty$, so $C^\varphi_1\ne0$ a.e.; the proof of Theorem~\ref{thm:tame}(c) then applies, with Jensen in place of Lemma~\ref{lem:trigpoly}.
\end{proof}

If V$(\gamma_1)$ holds for $\gamma_1=\gamma_1(\chi_{-4})=6.0209\ldots$, then at $M=\frac12$, with simple zeros, the constant in (b) at $c=0$ is $2.0549$ (numerical); e.g.\ $\nu_{1/2}((-\eps,\eps))\le0.0892$ at $\eps=10^{-10}$, which is weaker than the unconditional bound of Theorem~\ref{thm:cert} (the point of (b) is the decay as $\eps\to0$). TS implies V$(\gamma_0)$ for some $\gamma_0$, and V is strictly weaker: $\Gamma_0=\{b+nd:n\ge0\}$ with $b/d\notin\QQ$ satisfies V$(b)$ but not TS, not FD, and has no tame direction with a finite nonzero fibre (\file{density3/MAHLER.md}, Proposition~3.6, Example~3.7). FD does not imply V (Example~3.8 there, given with less detail). V for any single zero implies that two ordinates of $L(s,\chi)$ have an irrational ratio, the $L(s,\chi)$-analogue of item~1 of \cite{HKMN}, which is open; like TS and FD it cannot be verified by a finite computation.

\begin{theorem}[low vertices; GRH, NCZ; \rd{1 read}]\label{thm:LV}
Let $\mu_T$ be the smallest $\gamma\in\Gamma_T$ that is a vertex of $\mathrm{conv}(\pm\Gamma_T)$ (the largest element of $\Gamma_T$ always is). Then for every $M>0$ and $c\in\RR$,
\[
\nu_M(\{c\})\le2\liminf_{T\to\infty}\frac{\log^+\mu_T}{\log T}.
\]
\end{theorem}

\begin{proof}[Proof (outline; \file{density3/MAHLER.md}, \S4)]
With $\delta_T(c)=-\mathfrak m(F_{M,T}-c)$ and $\tau_T=\log(1/V_M^{>T})=\log T+O(\log\log T)$: if $\nu_M(\{c\})=\eta>0$, then on $S=\{F_M=c\}$ one has $F_{M,T}\circ\pi_T-c=-(F_M-F_{M,T}\circ\pi_T)$, and concavity of $\log$ on $S$ and $S^c$ gives $\delta_T(c)\ge\frac\eta2\tau_T-O(1)$ for every $T$. Lemma~\ref{lem:JN} at the vertex $\mu_T$ gives $\delta_T(c)\le\log(|M+i\mu_T|/2M)$.
\end{proof}

So ``low vertices'' ($\log\mu_T=o(\log T)$ along a sequence) exclude atoms. In rank one with a dense spectrum the only vertex is the top ordinate, and the bound is void; all the atom models of Section~\ref{sec:barriers} are of this kind.

\subsection{Graded linear independence}\label{sec:GLI}
For a finitely supported $k\in\ZZ^{(\Gamma)}$ the \emph{length} is $\|k\|_1$, and $k$ is a \emph{relation} if $\sum k_\gamma\gamma=0$. For $L\ge1$ let $S_L$ be the set of ordinates occurring in some nonzero relation of length $\le L$.

\begin{lemma}[moment comparison; no hypothesis; \rd{2 reads}]\label{lem:moment}
Let $N\ge1$ and $W,Z$ real random variables with $EZ^j=EW^j$ for $0\le j\le2N$. Suppose $W$ has a density, and $W=\mu+U+V$ with $|U|\le b$ a.s.\ and $E|V|^{2N}\le2\cdot N!\cdot(2\sigma^2)^N$. For $x\in\RR$ put $R=16(|x-\mu|+b+2\sqrt N\sigma)$ and $D_x=\operatorname{ess\,sup}_{|t-x|\le R}f_W(t)$. Then $P(|Z-x|\le2R/(N+1))\le352D_xR/(N+1)+56\cdot16^{-N}$.
\end{lemma}

\begin{proof}[Proof (outline; \file{density2/INVLO.md}, Lemma~A)]
Test against $p(t)=q((t-x)/R)^2$, where $q$ is a ``needle'' polynomial of degree $\le N$ built from the Fej\'er kernel: $q(0)=1$, $0\le q\le2$ on $[-1,1]$, $\int_{-1}^1q^2\le8\pi/(N+1)$, $q\ge0.27$ on $|s|\le2/(N+1)$, and $|q(s)|\le2(2|s|)^N$ for $|s|\ge1$. Then $P(|Z-x|\le2R/(N+1))\le14Ep(Z)=14Ep(W)$, and $Ep(W)$ is bounded using the density near $x$ and the moment bound far away.
\end{proof}

The rate $N^{-1/2}$ is sharp (Gauss--Hermite quadrature of a Gaussian). The following is a form of the inverse Littlewood--Offord principle (cf.\ \cite{NV}) on a subtorus: \emph{near-atoms force short integer relations}.

\begin{proposition}[finite models; no hypothesis; \rd{2 reads}]\label{prop:Bprime}
Let $\Gamma'\subset(0,\infty)$ be finite, $c_\gamma\in\CC$, $F=1+\re\sum_{\gamma\in\Gamma'}c_\gamma e_\gamma$ on the closure torus of $\{(e^{i\gamma y})_{\gamma\in\Gamma'}\}$, $L\ge2$, $N=\lfloor L/2\rfloor$, $B=\sum_{\gamma\in S_L}|c_\gamma|$, $\sigma_Y^2=\sum_{\gamma\in\Gamma'\setminus S_L}|c_\gamma|^2$, and let $D$ bound the density of $Y=\re\sum_{\gamma\in\Gamma'\setminus S_L}c_\gamma e^{i\theta_\gamma}$ with independent uniform phases. Then, with $R=16(|x-1|+B+2\sqrt N\sigma_Y)$, $\lambda\{|F-x|\le2R/(N+1)\}\le352DR/(N+1)+56\cdot16^{-N}$.
\end{proposition}

\begin{proof}
Let $W=1+X+Y$ with $X=\re\sum_{\gamma\in S_L}c_\gamma e_\gamma$ (Haar on its own torus) independent of $Y$. Expanding $(F-1)^j$ and $(W-1)^j$ into sequences of characters, a sequence contributes iff its net vector $k$ is a relation (resp.\ a relation supported on $S_L$); for $j\le L$ these agree, as nonzero relations of length $\le L$ are supported in $S_L$. So the first $L\ge2N$ moments agree, and Lemma~\ref{lem:moment} applies with $\mu=1$, $U=X$, $V=Y$ (Hoeffding).
\end{proof}

\begin{hypothesis}[GLI$(\psi)$]
$\psi:[0,\infty)\to[0,\infty)$ is nondecreasing, and every nonzero relation $k$ has $\|k\|_1\ge\psi(\max\{\gamma:k_\gamma\ne0\})$.
\end{hypothesis}

\begin{theorem}[GRH, NCZ, GLI$(\psi)$ with $\psi(t)\log(t+3)/t\to\infty$; \rd{2 reads}]\label{thm:GLI}
For every $M>0$ the measure $\nu_M$ has no atoms.
\end{theorem}

\begin{proof}[Proof (outline; \file{density2/INVLO.md}, Theorem~C)]
Let $T_L=\sup\{t:\psi(t)\le L\}$; then $S_L\subset\Gamma_{T_L}$, so $B\le C_\chi M\log^2(T_L+3)$ (here and below $C_\chi$, $C$ denote constants depending on $\chi$, resp.\ on $c,M,\chi$). The block $I=\Gamma\cap(T_L,2T_L]$ lies outside $S_L$ and contains $\gg T_L$ distinct ordinates (Riemann--von Mangoldt; zeros are not assumed simple), with coefficients between $M/(2T_L)$ and $4MC_\chi\log(2T_L+3)/T_L$; a density bound for random cosine sums (product of Bessel functions $J_0$) gives $D\le C(1+(T_L/\log T_L)^{1/2})$. Proposition~\ref{prop:Bprime} for the truncation $F_{M,T}$ (with $T\ge T_L$, $N=\lfloor L/2\rfloor$, $\rho=2R/(N+1)$), plus the truncation error $\le4V_M^{>T}/\rho^2$ and then $T\to\infty$, gives $\nu_M(\{c\})\le C(1+(T_L/\log T_L)^{1/2})[(1+\log^2T_L)/L+L^{-1/2}]+o(1)$, and $\psi(t)\ge Kt/\log t$ for large $t$ makes this $\ll K^{-1/2}$.
\end{proof}

\begin{corollary}[what an atom forces; GRH, NCZ; \rd{2 reads}]\label{cor:GLIatom}
If $\nu_M(\{c\})>0$, there is $C<\infty$ such that infinitely many nonzero relations $k$ satisfy $\|k\|_1\le C\,t_k/\log(t_k+3)$, where $t_k$ is the largest ordinate involved. More strongly (\file{density2/INVLO.md}, Proposition~E): if $\nu_M(\{c\})=\alpha$, then for every $\delta\in(0,1]$, every large $L$ and every $T$ with $\log^3L\le T\le C_2^{-1}\alpha^2\delta L\log L$, all but $\delta T\log T$ of the zeros in $(T,2T]$ (with multiplicity) lie in relations of length $\le L$.
\end{corollary}

LI implies GLI$(\psi)$ for every $\psi$, and GLI is incomparable with TS, FD and TS$_\ZZ$ (\file{density2/INVLO.md}, \S3; e.g.\ $\{\alpha u_n/2^{e_n}\}$ with distinct odd primes $u_n$ and rapidly increasing $e_n$ satisfies GLI and fails TS, FD, TS$_\ZZ$). Proposition~\ref{prop:Bprime} is effective, but its constants are far too large to improve Theorem~\ref{thm:cert} from a finite search for relations.

\subsection{Prime-side hypotheses}\label{sec:prime}
Fix $v\in C_c^\infty(-\frac12,\frac12)$ even, $v\ge0$, $\int v=1$, and $w=v*v$; then $\hat w=\hat v^2\ge0$ is even and Schwartz, $\hat w(0)=1$, $\hat w>0$ on $(-\pi,\pi)$; put $c_w=\min_{|\xi|\le1}\hat w(\xi)>0$. For $T>0$ the \emph{windowed prime sum} is $\Pi_T(y)=\sum_n\chi(n)\Lambda(n)n^{-1/2}\,Tw(T(y-\log n))$, a weighted sum over a window of relative width $\asymp1/T$ around $x=e^y$. By the explicit formula, $\Pi_T(y)=\Omega_T(y)-G_T(\iota(y))$ with $G_T=\sum_\gamma m_\gamma\hat w(\gamma/T)(e_\gamma+e_{-\gamma})$ and $\Omega_T(y)\to0$ (numerically confirmed for $\chi_{-4}$ to $5.4\cdot10^{-4}$). So the log-averaged moments $M_k(T)=\lim_Y\frac1Y\int^Y\Pi_T(y)^kdy$ exist and equal $(-1)^k\int G_T^k\,d\hA$. List the zeros with positive ordinate with multiplicity as $Z^+$, put $a_\rho=\hat w(\gamma_\rho/T)$, and call a nonzero $\kappa\in\ZZ^{(Z^+)}$ with $\sum\kappa_\rho\gamma_\rho=0$ a \emph{multiset relation} (two copies of a multiple zero give one of length $2$). Let $\mu_k(T)=E[(\sum_\rho2a_\rho\cos\theta_\rho)^k]$ with independent uniform phases.

\begin{theorem}[duality; GRH, NCZ; \rd{2 reads}]\label{thm:P1}
For all $T>0$, $k\ge1$: $(-1)^kM_k(T)=\mu_k(T)+\Delta_k(T)$, where $\Delta_k(T)=\sum'\prod_{i\le k}a_{\rho_i}\ge0$, the sum over sequences $(\rho_i,\sigma_i)\in(Z^+\times\{\pm1\})^k$ whose net vector is a multiset relation. In particular odd moments are $\le0$; for each $L$, $\Delta_\ell(T)=0$ for all $\ell\le L$ iff there is no multiset relation of length $\le L$ among the $\rho$ with $a_\rho>0$; each relation of length $\ell$ among zeros of height $\le T$ contributes $\ge c_w^\ell$ to $\Delta_\ell(T)$; and $\mu_k(T)$ is a polynomial in finite prime sums (explicit formula at $y=0$).
\end{theorem}

\begin{proof}
Expand $E_{\hA}[G_T^k]$; $E[e_g]=\mathbf 1_{g=0}$, while in the independent model a sequence contributes iff its net vector is $0$. The difference is the sum over sequences with nonzero net vector and total frequency $0$, all terms $\ge0$. The other statements follow by realising a relation of length $\ell$ by a sequence of length $\ell$, padding with cancelling pairs, and the explicit formula for $\sum_\rho\hat w(\gamma_\rho/T)^{2j}$.
\end{proof}

So at this level a prime-side correlation hypothesis is a zero-side relation hypothesis in disguise. Its quantitative form (moments up to order $L$ match the prediction to within $c_w^\ell$ at one scale $T\ge C_1L^{3/2}\log^2L$) bounds every atom by $C(1+|c-1|)L^{-1/2}$, via Proposition~\ref{prop:Bprime} (\file{density2/PRIMECORR.md}, Theorem~P2, \rd{2 reads}); the infinite version is equivalent to LI with simple zeros. Mixed moments of two windows with disjoint frequency bands do better.

\emph{Band decoupling.} Let $L\ge2$, $T^*\ge1$, a band $0<s_1<s_2<T^*$ and a margin $g\in(0,1)$. Take even windows $\beta_I,\beta_R\in C_c^\infty(\RR)$ with values in $[0,1]$, $\beta_I=1$ on $[s_1+g,s_2-g]$ and $\operatorname{supp}\beta_I\cap[0,\infty)\subset[s_1,s_2]$, $\beta_R=1$ on $[0,s_1-g]\cup[s_2+g,T^*]$ and $\operatorname{supp}\beta_R\cap[0,\infty)\subset[0,s_1]\cup[s_2,T^*+1]$. With $b=\check\beta$, the limits $\Pi_\beta(y)=\lim_X\sum_{n\le X}\chi(n)\Lambda(n)n^{-1/2}b(y-\log n)$ exist and $\Pi_\beta(y)=-G_\beta(\iota(y))+o(1)$, $G_\beta=\sum_\gamma m_\gamma\beta(\gamma)(e_\gamma+e_{-\gamma})$ (GRH, NCZ; \file{density2/PRIMECORR.md}, Lemma~P0$'$).

\begin{hypothesis}[H$_{\rm dec}(L;s_1,s_2,g,T^*)$, prime side]
\textup{(Z)} No zero has ordinate in $(s_1-g,s_1+g)\cup(s_2-g,s_2+g)$. \textup{(N)} The band contains a zero: $n_I:=\sum_{s_1<\gamma<s_2}m_\gamma\ge1$. \textup{(D)} For all $a\ge0$, $b\ge1$ with $a+b\le L$, the log-averaged joint moment $M_{a,b}=\lim_Y\frac1Y\int_1^Y\Pi_R^a\Pi_I^b\,dy$ satisfies $(-1)^{a+b}M_{a,b}-(-1)^aM_{a,0}\,\mu_b(n_I)<1$, where $\mu_b(n)=E[(\sum_{i\le n}2\cos\theta_i)^b]$.
\end{hypothesis}

Each quantity in H$_{\rm dec}$ is a limit of prime sums: (Z) and (N) are one-level statements via the explicit formula. Under (Z), the excess in (D) is a sum of nonnegative weights over sequences whose net vector is a relation involving the band, each relation among zeros of height $\le T^*$ having weight exactly $1$ (\file{density2/PRIMECORR.md}, Lemma~P1$'$); in particular (D) at $(a,b)=(0,2)$ forces the zeros in the band to be simple.

\begin{theorem}[GRH, NCZ, band decoupling; \rd{2 reads}]\label{thm:P7}
Suppose that along a sequence $L_j\to\infty$ the hypothesis H$_{\rm dec}(L_j;s_{1,j},s_{2,j},g_j,T^*_j)$ holds, including \textup{(N)}, with $T^*_j=\lfloor L_j/2\rfloor^4$, $2s_{1,j}\le s_{2,j}\le3s_{1,j}$ and $s_{1,j}/\log(s_{1,j}+3)=o(L_j)$. Then $\nu_M$ has no atoms, for every $M>0$.
\end{theorem}

\begin{proof}[Proof (outline; \file{density2/PRIMECORR.md}, \S3)]
If $s_{1,j}$ stays bounded along a subsequence, some zero $\gamma_0$ lies in infinitely many bands, and (D) with $L_j,T^*_j\to\infty$ excludes every relation involving $\gamma_0$; so $\gamma_0$ is simple and $\gamma_0\notin\spanQ{\Gamma\setminus\{\gamma_0\}}$, TS holds, and Theorem~\ref{thm:TS} applies. If $s_{1,j}\to\infty$: by Lemma~P1$'$ and (D) no relation of length $\le L$ among zeros of height $\le T^*$ involves the band, so the band, which contains $\gg s_1$ distinct ordinates, is free; Proposition~\ref{prop:Bprime} with the density bound $D\le C(1+(s_1/\log s_1)^{1/2})$ and the truncation at $T^*=N^4$ give $\nu_M(\{c\})\le C(1+(s_1/\log s_1)^{1/2})(\log^2L/L+L^{-1/2})+o(1)\to0$.
\end{proof}

The first read of this theorem (\file{density2/REFEREE-4b.md}) found that its first version, without (N), held vacuously for an empty band and was therefore false as stated; condition (N) and the treatment of bounded bands are the corrections it proposed. A second, independent read (\file{density2/REFEREE-5.md}) found the corrected theorem, in both cases, correct as stated. LI with simple zeros implies H$_{\rm dec}$ for bands at every large height. For spectra with simple zeros (in one admissible band per scale), GLI implies H$_{\rm dec}$, so Theorem~\ref{thm:GLI} restricted to such spectra is a special case of Theorem~\ref{thm:P7}; without simplicity it is not (a $\QQ$-independent spectrum with all multiplicities $2$ satisfies GLI but fails (D)). Devin's Theorem~5.4(1) does not imply H$_{\rm dec}$ (\file{density2/PRIMECORR.md}, \S3). (The example offered in the working note to show that H$_{\rm dec}$ is strictly weaker than GLI, and that it does not imply Devin's Theorem~5.4(1), fails condition (D), as the second read found; we make neither claim.)

\begin{remark}[heuristic; not a theorem]
Pointwise Hardy--Littlewood asymptotics for $k$-tuples of shifted primes with power-saving error terms, as used by Montgomery and Soundararajan \cite{MS}, do not reach these moments: for windows of fixed relative width $1/T$ the error terms exceed the main term by a factor $\asymp(x/T)^{k/2}x^{-1/2+\eps}$. And by Corollary~\ref{cor:GLIatom} the relations an atom needs may all be far longer than any fixed moment order, so a hypothesis on moments of fixed order $K$ (or $K=o(T/\log T)$), with any precision, cannot exclude them.
\end{remark}
